\begin{exercise}[Point-set definition]\label{exr:vi35-01}
For \(S=k[x,y,z]\), explain why \((x,y,z)\) is not a point of \(\Proj S\).
\end{exercise}

\begin{exercise}[Homogeneous prime]\label{exr:vi35-02}
Show \((x)\subset k[x,y]\) defines a point of \(\Proj k[x,y]\).
\end{exercise}

\begin{exercise}[Closed-set union]\label{exr:vi35-03}
Prove \(V_+(I)\cup V_+(J)=V_+(IJ)\) for homogeneous ideals.
\end{exercise}

\begin{exercise}[Closed-set intersection]\label{exr:vi35-04}
Prove \(\bigcap_\alpha V_+(I_\alpha)=V_+(\sum I_\alpha)\).
\end{exercise}

\begin{exercise}[Standard-open intersection]\label{exr:vi35-05}
Prove \(D_+(f)\cap D_+(g)=D_+(fg)\).
\end{exercise}

\begin{exercise}[Standard cover]\label{exr:vi35-06}
If \(S_+=(f_0,\ldots,f_r)\) with homogeneous \(f_i\), prove \(\Proj S=\bigcup_iD_+(f_i)\).
\end{exercise}

\begin{exercise}[One-variable Proj]\label{exr:vi35-07}
Compute \(\Proj k[x]\) and its structure sheaf.
\end{exercise}

\begin{exercise}[Empty Proj]\label{exr:vi35-08}
Compute \(\Proj k[\varepsilon]/(\varepsilon^3)\) when \(\deg\varepsilon=1\).
\end{exercise}

\begin{exercise}[Affine chart]\label{exr:vi35-09}
Compute \(k[x,y,z]_{(x)}\) in the standard grading.
\end{exercise}

\begin{exercise}[Overlap ring]\label{exr:vi35-10}
For \(S=k[x,y]\), compute the coordinate ring of \(D_+(x)\cap D_+(y)\).
\end{exercise}

\begin{exercise}[Homogeneous quotient]\label{exr:vi35-11}
Show set-theoretically \(\Proj(S/I)\cong V_+(I)\) for homogeneous \(I\).
\end{exercise}

\begin{exercise}[Stalk fractions]\label{exr:vi35-12}
Describe elements of \(S_{(\mathfrak p)}\) as fractions.
\end{exercise}

\begin{exercise}[Locality]\label{exr:vi35-13}
Prove the fractions with numerator in \(\mathfrak p\) form the maximal ideal of \(S_{(\mathfrak p)}\).
\end{exercise}

\begin{exercise}[Associated sheaf]\label{exr:vi35-14}
For a graded module \(M\), state the sections of \(\widetilde M\) on \(D_+(f)\).
\end{exercise}

\begin{exercise}[Stalk of a graded sheaf]\label{exr:vi35-15}
Show \(\widetilde M_{\mathfrak p}\cong M_{(\mathfrak p)}\).
\end{exercise}

\begin{exercise}[Structure sheaf]\label{exr:vi35-16}
Prove \(\widetilde S\cong\OO_{\Proj S}\).
\end{exercise}

\begin{exercise}[Exactness]\label{exr:vi35-17}
Explain why the functor \(M\mapsto\widetilde M\) is exact on degree-zero exact sequences of graded modules.
\end{exercise}

\begin{exercise}[Shift notation]\label{exr:vi35-18}
Write the chart sections of \(\OO(d)=\widetilde{S(d)}\).
\end{exercise}

\begin{exercise}[Reducible example]\label{exr:vi35-19}
Compute \(\Proj k[x,y]/(xy)\).
\end{exercise}

\begin{exercise}[Nonreduced example]\label{exr:vi35-20}
Compute the x-chart of \(\Proj k[x,y]/(y^2)\).
\end{exercise}

\begin{exercise}[Veronese prototype]\label{exr:vi35-21}
For \(S=k[x,y]\), compare the x-chart of \(\Proj S\) with the \(x^2\)-chart of \(\Proj S^{(2)}\).
\end{exercise}

\begin{exercise}[Nonhomogeneous warning]\label{exr:vi35-22}
Explain why \(x+y^2\) is not a valid homogeneous projective equation in standard \(k[x,y]\).
\end{exercise}

\begin{exercise}[Subspace topology]\label{exr:vi35-23}
Show the \(V_+\)-topology agrees with the subspace topology from \(\Spec S\).
\end{exercise}

\begin{exercise}[Boundary to VI/36]\label{exr:vi35-24}
For \(S=k[x_0,\ldots,x_n]\), state without developing projective-space theory what the standard charts are.
\end{exercise}